\section{Introduction}

The cost of solving a convex optimization problem depends on both its
mathematical iteration and its implementation. A method requiring fewer
iterations can have a more expensive update; lower precision can reduce
arithmetic cost while limiting attainable accuracy; and evaluating a stopping
criterion can consume a substantial part of execution time. Practical solvers
therefore combine algorithmic choices with preprocessing, scheduling and
implementation decisions \citep{Applegate2021Practical,Stellato2020OSQP}.
For repeated solves from a problem distribution on a fixed device, the
relevant objective is time to a solution meeting the declared numerical
criterion, together with the cost of finding and preparing the solver.

\textbf{OptEvolve} represents this decision as search over executable solver
recipes. A recipe specifies a base operator splitting, parameters and supported
wrappers, together with precision, certificate cadence and an eligible kernel
implementation. A construction gate checks sufficient convergence conditions
before the numerical iteration is executed. The compiler assembles the
iteration from implemented atoms. Numerical termination is checked separately
on the original problem, including when the iteration uses lower precision
or a rescaled problem.

LLM proposals, evolutionary populations and execution feedback have substantial
prior art \citep{RomeraParedes2023Mathematical,Ye2024ReEvo,Agrawal2025GEPA}.
Automated optimization research also includes numerical method design,
symbolic proof discovery and formal verification \citep{Kim2026AutoOPT}.
Our question concerns their use within a restricted executable recipe space:
does hardware-aware selection improve time to a checked solution after strong
algorithm and implementation tuning, and when must these choices be made
jointly? We treat an advantage of joint search or LLM guidance as an empirical
question, not a consequence of including those components.

\paragraph{Scope and contributions.}
The current system provides a typed representation of implemented splitting
methods with a construction theorem for base schemes and relaxation
(Section~\ref{sec:theory}); a hardware-aware execution layer with numerical
admission and original-problem checks (Section~\ref{sec:hardware}); and a
versioned evaluation protocol that separates construction admission, successful
termination, search cost and execution cost. Preliminary TV and LP experiments
identify measurable execution-policy gains over legacy execution, failures of
pure float32 at tight tolerance, and a finite LP catalog in which joint and
sequential selection coincide (Section~\ref{sec:experiments}). These diagnostics
motivate controlled comparisons against stronger tuned baselines. They do not
yet demonstrate a new optimization algorithm, an advantage from generated
kernels, or better search efficiency from an LLM.
